\documentclass[11pt,a4paper,oneside]{article}

% --- Pacotes de Idioma e Codificação ---
\usepackage[utf8]{inputenc}
\usepackage[portuguese]{babel}
\usepackage[T1]{fontenc}

% --- Pacotes de Design e Layout ---
\usepackage{geometry}
\geometry{
    a4paper,
    total={170mm,257mm},
    left=25mm,
    top=25mm,
}
\usepackage{xcolor}
\usepackage{titlesec}
\usepackage{enumitem}
\usepackage{booktabs}
\usepackage{fancyhdr}
\usepackage{hyperref}
\usepackage{tcolorbox}
\usepackage{tabularx}

% --- Definição de Cores (Paleta Profissional Médica) ---
\definecolor{MedBlue}{RGB}{0, 51, 102}
\definecolor{SoftBlue}{RGB}{230, 242, 255}
\definecolor{DarkGrey}{RGB}{50, 50, 50}
\definecolor{AlertRed}{RGB}{180, 0, 0}

% --- Configuração de Fontes e Estilos ---
\titleformat{\section}{\color{MedBlue}\normalfont\Large\bfseries}{\thesection}{1em}{}
\titleformat{\subsection}{\color{MedBlue}\normalfont\large\bfseries}{\thesubsection}{1em}{}
\renewcommand{\familydefault}{\sfdefault}

% --- Cabeçalho e Rodapé ---
\pagestyle{fancy}
\fancyhf{}
\lhead{\footnotesize \textcolor{MedBlue}{Dossiê Estratégico IntraClinica}}
\rhead{\footnotesize \textcolor{DarkGrey}{Confidencial | Fevereiro 2026}}
\cfoot{\thepage}

% --- Configurações de Links ---
\hypersetup{
    colorlinks=true,
    linkcolor=MedBlue,
    filecolor=magenta,      
    urlcolor=cyan,
    pdftitle={Dossiê Estratégico IntraClinica 2026},
}

\begin{document}

% --- Capa ---
\begin{titlepage}
    \centering
    \vspace*{2cm}
    {\Huge \bfseries \textcolor{MedBlue}{IntraClinica}} \\
    \vspace{0.5cm}
    {\Large \textcolor{DarkGrey}{Soberania Digital e Inteligência Clínica Aplicada}} \\
    \vspace{2cm}
    
    \begin{tcolorbox}[colback=SoftBlue, colframe=MedBlue, arc=5pt, width=0.9\textwidth, center]
        \centering
        \textbf{\large Relatório de Governança, Expansão e Visão de Produto} \\
        \vspace{0.2cm}
        \textit{Documento Técnico-Estratégico para Diretoria e Curadoria Médica}
    \end{tcolorbox}
    
    \vfill
    
    \begin{minipage}{0.4\textwidth}
        \begin{flushleft} \large
            \textbf{Preparado por:} \\
            Bruno \\
            Axi (Digital Peer)
        \end{flushleft}
    \end{minipage}
    \begin{minipage}{0.4\textwidth}
        \begin{flushright} \large
            \textbf{Status:} Ativo \\
            \textbf{Versão:} 2.1.0 \\
            \textbf{Data:} 20/02/2026
        \end{flushright}
    \end{minipage}
    
    \vspace{1cm}
\end{titlepage}

\newpage
\tableofcontents
\newpage

\section{Introdução e Visão Estratégica}
O projeto \textbf{IntraClinica} foi fundamentado na premissa de que a tecnologia médica atual sofre de uma "paralisia por complexidade". Softwares tradicionais oferecem módulos excessivos, mas falham em resolver as dores viscerais do cotidiano médico: a perda de controle sobre o estoque e a burocratização do atendimento.

Nossa visão é pautada pelo binômio \textbf{``Medicine for Humans. Intelligence for Business.''}. Buscamos devolver a soberania ao médico através de uma infraestrutura que ele controla fisicamente e logicamente.

\section{Diagnóstico Clínico de Campo}
A concepção desta arquitetura nasceu da observação direta de gargalos críticos em unidades de saúde de médio porte:

\begin{enumerate}
    \item \textbf{Hemorragia Financeira no Almoxarifado:} Insumos como kits de sutura e medicamentos desaparecem sem rastro devido a controles manuais falhos. O impacto é uma perda direta de margem que muitas vezes a clínica nem consegue quantificar.
    \item \textbf{Fadiga de UX (Síndrome do Sistema Inchado):} Médicos e enfermeiros evitam o uso de softwares legados por serem hostis ao fluxo de trabalho rápido. O resultado é o preenchimento de dados apenas "para o faturamento", perdendo-se a riqueza do dado clínico.
\end{enumerate}

\section{Soluções e Use-Cases de Valor (Experiência Clínica Soberana (Ecossistema IntraClinica))}

\subsection{O Guardião de Estoque Inteligente}
O IntraClinica transforma o estoque de uma planilha passiva para um \textbf{agente ativo de auditoria}.
\begin{itemize}
    \item \textbf{Fluxo de Baixa Automática:} Ao finalizar um procedimento no prontuário (ex: uma sutura simples), o sistema consulta a "receita" do procedimento e exige a confirmação da baixa dos itens associados (fio nylon, kit sutura, antisséptico).
    \item \textbf{Auditoria por IA:} Modelos de inteligência local cruzam o volume de procedimentos realizados com a flutuação do estoque, emitindo alertas de anomalia se o consumo divergir drasticamente do esperado para aquele turno.
\end{itemize}

\subsection{Atendimento com a ``Regra de Uma Mão''}
Reconhecendo que o médico frequentemente utiliza tablets ou dispositivos móveis durante a anamnese, a interface foi desenhada para operação monomanual.
\begin{itemize}
    \item \textbf{Design Sem Fricção:} Navegação por \textit{Signals} do Angular 21, permitindo atualizações de estado instantâneas e transições fluidas entre recepção, agendamento e prontuário.
    \item \textbf{Foco no Paciente:} O objetivo é que o registro clínico seja um assistente de fundo, preservando o contato visual e a empatia, pilares da boa prática médica.
\end{itemize}

\subsection{Engajamento e Comunicação Social}
O IntraClinica inclui um módulo de inteligência criativa focado no fortalecimento da presença digital da clínica.
\begin{itemize}
    \item \textbf{Criação de Conteúdo:} O sistema é capaz de gerar rascunhos de postagens para redes sociais baseados em temas clínicos ou de saúde preventiva.
    \item \textbf{Automação Ética:} A comunicação é desenhada para educar o paciente e reforçar a autoridade médica, mantendo a sobriedade exigida pelo conselho de classe.
\end{itemize}

\section{Arquitetura e Soberania de Dados}
Diferente das soluções tradicionais em nuvem opaca, o IntraClinica utiliza uma abordagem \textbf{Soberana}, baseada no modelo relacional normalizado detalhado no diagrama \texttt{db\_schema.svg}.

\begin{itemize}
    \item \textbf{Backend-as-a-Service Resiliente:} Implementado sobre Supabase (PostgreSQL 16), garantindo que a clínica seja dona da base de dados e possa realizar auditorias granulares.
    \item \textbf{Isolamento Multi-Tenancy (RLS):} Utilização de \textit{Row Level Security} nativo do Postgres para garantir o isolamento lógico rigoroso dos dados entre diferentes unidades. Um médico só acessa dados onde o identificador da clínica coincide com seu perfil autorizado.
    \item \textbf{Máquina de Estados de Agendamento:} Lógica rigorosa de transição de status (Agendado $\rightarrow$ Aguardando $\rightarrow$ Em Atendimento $\rightarrow$ Realizado), impedindo saltos de etapa e garantindo a integridade do faturamento e dos indicadores de produtividade.
    \item \textbf{Segurança Zero-Trust:} Cada requisição de dado é validada no nível do banco de dados, protegendo a rede contra acessos indevidos.
\end{itemize}

\section{Governança e Matriz de Acesso (IAM)}
A expansão para novos médicos parceiros é protegida por uma estrutura de \textbf{Identity and Access Management} rigorosa:

\begin{table}[h!]
\centering
\caption{Matriz de Funções e Permissões}
\vspace{0.3cm}
\begin{tabularx}{\textwidth}{@{}lXl@{}}
\toprule
\textbf{Papel (Role)} & \textbf{Descrição do Escopo} & \textbf{Permissão Core} \\ \midrule
\textcolor{MedBlue}{Diretoria} & Gestão de protocolos, curadoria e auditoria global. & \texttt{admin.full} \\
\textcolor{MedBlue}{Médico} & Atendimento pleno e gestão de prontuário soberano. & \texttt{clinical.write} \\
\textcolor{MedBlue}{Estoque} & Gestão total de suprimentos e custos de aquisição. & \texttt{inventory.manage} \\
\textcolor{MedBlue}{Recepção} & Foco em agenda, check-in e cadastro demográfico. & \texttt{patient.read} \\ \bottomrule
\end{tabularx}
\end{table}

\section{Roadmap de Evolução Cognitiva}
Estamos em fase de transição de uma malha de dados puramente geométrica para uma **Topologia Neural (Tier 4)** via Axio Nexus:
\begin{description}
    \item[Fase 1 (Atual):] Consolidação da infraestrutura Supabase e automação de fluxos de estoque.
    \item[Fase 2 (Próxima):] Implementação de assistentes de anamnese que "escutam" a consulta e geram rascunhos de evolução, aguardando apenas a assinatura do médico.
    \item[Fase 3 (Expansão):] Interconexão de unidades parceiras via protocolo de memória distribuída, permitindo curadoria de casos clínicos de forma anonimizada e soberana.
\end{description}

\section{Conclusão}
O IntraClinica posiciona-se como a espinha dorsal de uma nova era na gestão médica — onde a tecnologia respeita a biologia e o lucro é uma consequência direta da eficiência operacional e da soberania do dado.

\vfill
\centering
\textit{Documento gerado e verificado conforme protocolos de engenharia v3.3.} 🧪🧬🏥
\end{document}
